\subsection{The auxiliary torus cover on exact finite arithmetic fibres}
\label{sec:ns-torus-finite}

The source of the operator in this subsection is the public manuscript
\emph{Finite Time Blowup for Navier--Stokes}, whose displayed author is
OpenAI, Section~6, equation~(6.2) on printed page~63 and Lemma~6.2,
equation~(6.19), on printed page~67.  The preserved source is the PDF
with SHA-256
\texttt{8c8a94ad9ac824c8b605b9827cadf7beaca}\allowbreak%
\texttt{48bd10b380de3cfc872a2c37afa81};
the associated source pin is
\texttt{8937a8f4cbc7abaab5e9e97d1cc7f5d2319d9538}.
The following proofs concern this explicitly specified integer operator
and its arithmetic counting consequences. They do not use a
Navier--Stokes existence or blowup theorem. Authorship of other results
is not inferred from this operator calculation.

Retain, in the coordinates and order of the source,
\[
 \begin{gathered}
 J=\begin{pmatrix}3&1\\1&5\end{pmatrix},\qquad \det J=14,\\
 v_r=(1,1-\sqrt2),\qquad v_t=(\sqrt2-1,1),\\
 Jv_r=(4-\sqrt2)v_r,\qquad Jv_t=(4+\sqrt2)v_t.
 \end{gathered}
\]
For example the two components of $Jv_r$ are $4-\sqrt2$ and
$6-5\sqrt2=(4-\sqrt2)(1-\sqrt2)$; those of $Jv_t$ are
$3\sqrt2-2=(4+\sqrt2)(\sqrt2-1)$ and $4+\sqrt2$.
No eigenvector is divided by its length. Put
$\mathbb T=\mathbb R/\mathbb Z$ and
$\Phi:\mathbb T^2\to\mathbb T^2$, $\Phi(x,y)=(3x+y,x+5y)$.

\begin{theorem}[Directional operators in the unchanged source coordinates]
\label{thm:ns-directional-intertwining}
Let $\mathcal P$ be the complex trigonometric polynomials on
$\mathbb T^2$, and let $\mathcal P_0$ consist of those with zero
constant Fourier coefficient. For $m\geq0$ define
$P_m:\mathcal P\to\mathcal P$ by $(P_m f)(Y)=f(J^mY)$.
Writing $f(Y)=\sum_{n\in\mathbb Z^2}\widehat f(n)
e^{2\pi i n\cdot Y}$ with finite support gives the exact formula
\[
 P_m f(Y)=\sum_n\widehat f(n)e^{2\pi i((J^m)^Tn)\cdot Y}.
                                                               \tag{T0a}
\]
Thus $P_m$ is injective; its image is precisely the polynomials
whose frequencies belong to $(J^m)^T\mathbb Z^2$, and its inverse
on that image has coefficient at $n$ equal to the coefficient at
$(J^m)^Tn$ of the image polynomial. It preserves the constant
coefficient and the probability Haar average.

For $(v,\lambda)=(v_r,4-\sqrt2)$ or $(v_t,4+\sqrt2)$, define
$D_v=v\cdot\partial_Y$. Then
\[
 D_vP_m=\lambda^mP_mD_v,\qquad
 D_v^{-1}P_m=\lambda^{-m}P_mD_v^{-1}\quad\text{on }\mathcal P_0,
                                                               \tag{T0b}
\]
where the exact inverse with its original Fourier constant is
\[
 D_v^{-1}f(Y)=\sum_{n\ne0}
       \frac{\widehat f(n)}{2\pi i(v\cdot n)}e^{2\pi i n\cdot Y}.
                                                               \tag{T0c}
\]
The kernel of $D_v:\mathcal P\to\mathcal P$ is exactly the
constant polynomials; its image is $\mathcal P_0$; and every
fibre at $f\in\mathcal P_0$ is $D_v^{-1}f+\mathbb C$.
\end{theorem}
\begin{proof}
The equality $n\cdot J^mY=((J^m)^Tn)\cdot Y$ proves (T0a).
The integer frequency map is injective because $\det J^m=14^m$
is nonzero. This proves the image, inverse and singleton fibres
of $P_m$, and the zero frequency is preserved in both directions.
Integrating each character over $[0,1]^2$ proves the averaging
claim directly. By the chain rule
$v\cdot\partial_Y f(J^mY)=(J^mv)\cdot(\partial f)(J^mY)$;
the displayed eigenvector computations give $J^mv=\lambda^mv$.
This proves the first identity (T0b) with the actual source vectors.
For nonzero $n=(n_1,n_2)\in\mathbb Z^2$, neither
$n_1+(1-\sqrt2)n_2$ nor $(\sqrt2-1)n_1+n_2$ is zero.
Otherwise, unless both integer coordinates were zero, solving
that equation would make $\sqrt2$ rational. To recall the exact
obstruction, a reduced fraction $a/b$ with square two satisfies
$a^2=2b^2$, forcing $a$ and then $b$ even, a contradiction.
Each nonconstant Fourier mode of $D_v$ therefore has nonzero
multiplier $2\pi i(v\cdot n)$. This proves (T0c), the kernel,
the image and every asserted fibre. On a pulled-back mode the
denominator is
$2\pi i(v\cdot(J^m)^Tn)=\lambda^m2\pi i(v\cdot n)$.
Its reciprocal proves the second identity (T0b).
\end{proof}

\begin{theorem}[The cover and its complete finite-torsion restriction]
\label{thm:ns-torus-fibres}
For a target $(a,b)\in\mathbb T^2$, choose real representatives
$\widetilde a,\widetilde b$. Its complete fibre under $\Phi$ is
\[
 \left\{\left(
 \frac{5\widetilde a-\widetilde b+k}{14},
 \widetilde a-3\frac{5\widetilde a-\widetilde b+k}{14}
 \right)\pmod{\mathbb Z^2}:k=0,\ldots,13\right\}.       \tag{T1}
\]
In particular $\Phi$ is surjective and
$\ker\Phi=\{(k/14,-3k/14):k=0,\ldots,13\}$.

For an integer $N\geq1$, identify $\mathbb T[N]$ with
$\mathbb Z/N\mathbb Z$ by $j\mapsto j/N+\mathbb Z$, with inverse
$z\mapsto N\widetilde z\pmod N$ for any real representative
$\widetilde z$ of $z$. Changing that representative by an integer
changes $N\widetilde z$ by a multiple of $N$.
The restriction $J_N$ to
$(\mathbb Z/N\mathbb Z)^2$ has image and kernel
\[
 \begin{split}
 I_N&=\{(a,b):g\mid5a-b\},\qquad g=\gcd(N,14),\\
 K_N&=\{(kN/g,-3kN/g)\pmod N:0\leq k<g\}.
 \end{split}                                                    \tag{T2}
\]
Divisibility in (T2) is independent of the chosen integer target
representatives. A target outside $I_N$ has empty fibre. For a target
inside $I_N$, put $n=N/g$. If $n=1$ put $x_0=0$; otherwise put
\[
 x_0\equiv (14/g)^{-1}(5a-b)/g\pmod n.
\]
Its entire fibre is
\[
 \{(x_0+kn,a-3x_0-3kn)\pmod N:0\leq k<g\}.           \tag{T3}
\]
Consequently $|K_N|=g$ and $|I_N|=N^2/g$.
\end{theorem}
\begin{proof}
The first equation $3x+y=a$ forces $y=a-3x$ in either group.
Substitution in the second equation gives $14x=5a-b$.
On $\mathbb T$ this equation has exactly the fourteen solutions in
(T1): after choosing real representatives the difference from
$(5\widetilde a-\widetilde b)/14$ belongs to
$\frac1{14}\mathbb Z/\mathbb Z$. Their first coordinates are distinct.
Changing representatives leaves the solution equations unchanged, so
leaves the displayed full fibre unchanged. On $\mathbb Z/N\mathbb Z$
the equation is solvable precisely when $g\mid5a-b$: necessity follows
because $g$ divides both $14$ and $N$; sufficiency follows by dividing
the integer congruence by $g$, since $\gcd(14/g,n)=1$.
The resulting single class modulo $n$ has exactly the $g$ lifts in
(T3). For $n=1$ that divided congruence has its single zero class.
At the zero target this gives (T2)'s kernel. Every occupied fibre has
$g$ elements and the domain has $N^2$ elements, proving the image
cardinality. The coordinate identification with torsion commutes with
$J$ because every matrix entry is an integer.
\end{proof}

\begin{theorem}[Exact averaging, aliasing, and the complete-preimage repair]
\label{thm:ns-torus-average}
For every function $f:(\mathbb Z/N\mathbb Z)^2\to\mathbb C$,
\[
 \frac1{N^2}\sum_{x,y\bmod N} f(3x+y,x+5y)
   =\frac g{N^2}\sum_{(a,b)\in I_N}f(a,b).             \tag{T4}
\]
This equals the uniform mean of $f$ on all targets for every $f$ if
and only if $g=1$. For a frequency $(u,v)\bmod N$, the character
$\exp(2\pi i(ua+vb)/N)$ pulls back to frequency
$(3u+v,u+5v)$. Its mean after pullback is one precisely at
\[
 (u,v)=(kN/g,-3kN/g)\pmod N,\qquad 0\leq k<g,          \tag{T5}
\]
and is zero at every other frequency.

For any finite subgroup $D\subset\mathbb T^2$, let
$H=\Phi^{-1}(D)$. Then $H$ is a finite subgroup,
$|H|=14|D|$, and $\Phi:H\to D$ has the complete fibres (T1).
Define typed linear maps
\[
 \begin{aligned}
 P:\mathbb C^D&\longrightarrow\mathbb C^H,& Pf(h)&=f(\Phi h),\\
 L:\mathbb C^H&\longrightarrow\mathbb C^D,& Lv(d)&=
       \frac1{14}\sum_{\Phi h=d}v(h).
 \end{aligned}                                                   \tag{T6}
\]
Then $LP$ is the identity, $PL$ is the projection onto functions
constant on every deck orbit $h+\ker\Phi$, and
\[
 \frac1{|H|}\sum_{h\in H}Pf(h)=\frac1{|D|}\sum_{d\in D}f(d),
 \qquad
 \langle Pf,v\rangle_H=\langle f,Lv\rangle_D.          \tag{T7}
\]
Here $\langle f,v\rangle_D=|D|^{-1}\sum f\overline v$,
and likewise on $H$. The full fibre of $L$ at $f$ is
$Pf+\ker L$, where
$\ker L=\{v:\sum_{\Phi h=d}v(h)=0\text{ for every }d\}$.
The map $P$ is injective, and its fibre at $v$ is the singleton
$\{Lv\}$ if $v$ is deck-invariant, and is empty otherwise.
\end{theorem}
\begin{proof}
Sum $f$ by the fibres (T3) to obtain (T4). If $g=1$ the image is the
whole group. If $g>1$, the function indicating the zero target has
pullback mean $g/N^2$ and full target mean $1/N^2$, proving failure
without any regularity assumption. For characters, the geometric
sum $\sum_{j=0}^{N-1}\exp(2\pi i kj/N)$ is $N$ if $N\mid k$
and zero otherwise: multiplication by its ratio leaves the sum
unchanged, so a ratio different from one forces it to vanish. Apply this to
both input coordinates and use (T2), since $J$ is symmetric.

The inverse image $H$ is a subgroup because $\Phi$ is a homomorphism.
It is the disjoint union of the fourteen-element fibres over $D$.
On every such fibre $Pf$ is the constant $f(d)$, which proves
$LP=1$ and the first identity (T7). Formula (T6) makes $PLv$ the
arithmetic mean of $v$ on its deck orbit. It therefore has image
exactly the deck-invariant functions and squares to itself.
Summing $f(d)\overline{v(h)}$ on these same fibres proves the
inner-product identity with the stated factors $|H|=14|D|$.
The kernel and all claimed fibres follow directly: $L(v-Pf)=0$
is equivalent to $Lv=f$, while $Pf=v$ implies and, for
deck-invariant $v$, is implied by $f=Lv$.
\end{proof}

For the specific subgroup $D=\mathbb T[N]^2$, (T1) gives points
of $H$ with denominators dividing $14N$, and $|H|=14N^2$.
Thus (T7) changes the sampling domain to this explicitly given
preimage; it does not assert that the original $N^2$ samples acquire
new values. The continuous identity in source equation~(6.19)
has the same Fourier calculation: on $\mathbb T^2$ a frequency
$n\in\mathbb Z^2$ has nonzero mean exactly when $n=0$, and
$J^Tn=0$ exactly when $n=0$ because $\det J=14$.
For completeness this identity holds for all continuous functions.
Set $K_M(t)=M^{-1}|\sum_{j=0}^{M-1}e^{2\pi ijt}|^2$.
Expansion shows that this is a nonnegative trigonometric polynomial;
integrating its expansion shows its mass is one, since exactly the
$M$ diagonal terms survive. Convolution with
$K_M(t_1)K_M(t_2)$ is a trigonometric polynomial and converges
uniformly to each continuous function. Indeed, outside any fixed
neighborhood of zero each one-dimensional kernel is bounded by
$1/(M\sin^2(\pi\delta))$; uniform continuity bounds the integral
inside that neighborhood and this bound controls the complement.
The finite Fourier identity therefore passes to the integral limit.
This argument uses the original torus and its probability Haar
measure, not a replacement of a finite exponent box by a subgroup.

\begin{proposition}[Powers and product covers, with all lift labels]
\label{prop:ns-torus-powers}
For $m\geq0$ the map $\Phi_m=J^m$ has torus fibres of size $14^m$.
For a target $d$ with real representative $\widetilde d$, the exact
parametrization is
\[
 \mathbb Z^2/J^m\mathbb Z^2\longrightarrow\Phi_m^{-1}(d),
 \qquad [k]\longmapsto J^{-m}(\widetilde d+k)\pmod{\mathbb Z^2}.
                                                               \tag{T8}
\]
Its inverse sends a real lift $\widetilde h$ to
$[J^m\widetilde h-\widetilde d]$. In particular a full preimage of
a finite subgroup $D$ has $14^m|D|$ elements, and (T6)--(T7)
hold with $\Phi_m$ and $14^m$. For a block diagonal product of $s$
copies of $J^m$, use $s$ independent copies of (T8); every fibre
has $14^{ms}$ elements and the averaging statements have that factor.

On fixed $N$-torsion, define $S_0(d)=\{d\}$ and
$S_{j+1}(d)=\coprod_{w\in S_j(d)}J_N^{-1}(w)$, with every
one-step fibre given by (T3). Then $S_j(d)=\{h:J_N^jh=d\}$
with no repeated elements. These exact recursion fibres must be
used in place of an assumed size $\gcd(N,14)^j$.
\end{proposition}
\begin{proof}
Changing $k$ by $J^m z$ changes (T8) by $z$ and hence changes no
torus point. Conversely equality of torus points implies this exact
change in $k$. Every lift in the fibre gives an integral difference
$J^m\widetilde h-\widetilde d$, proving surjectivity and the inverse.
One-step surjectivity and fourteen-element fibres imply by induction
that the fibre of $\Phi_m$ consists of $14^m$ points: the preimages
of distinct intermediate points are disjoint. This proves the
cardinality of the quotient as well. Fibre counting now proves all
averaging assertions exactly as in (T6)--(T7), and independent
coordinate choices prove the product statement. In the torsion
recursion each candidate $h$ has the unique intermediate point
$J_Nh$, so the displayed union is disjoint and the induction gives
precisely $J_N^{j+1}h=d$. For example, at $N=2$, $J_N$ is the
matrix whose four entries are one; its square is zero. Its kernel
has two elements at $j=1$ and four elements at every $j\geq2$,
which also verifies why the proposed exponential size is incorrect.
\end{proof}

\begin{lemma}[The actual finite dual and its coordinate embedding]
\label{lem:ns-finite-dual}
Let $G$ be a finite abelian multiplicative group and
$\widehat G=\operatorname{Hom}(G,\{z\in\mathbb C:|z|=1\})$.
Then $|\widehat G|=|G|$, its characters separate the elements of
$G$, and
\[
 \frac1{|G|}\sum_{\chi\in\widehat G}\chi(u)
      =\begin{cases}1&u=1,\\0&u\ne1.\end{cases}       \tag{T9}
\]
Choose and retain an ordered generating tuple $g_1,\ldots,g_s$ of
$G$ (the list of all elements is permitted). The map
\[
 E:\widehat G\to\mathbb T^s,\qquad
 E(\chi)_j=\theta_j\pmod{\mathbb Z},\quad
       e^{2\pi i\theta_j}=\chi(g_j)                  \tag{T10}
\]
is an injective homomorphism. Its inverse on the image sends
$\theta$ to the character
$g_1^{n_1}\cdots g_s^{n_s}\mapsto
\exp(2\pi i\sum n_j\theta_j)$. The image consists exactly of
the vectors for which $\sum n_j\theta_j=0$ in $\mathbb T$
for every relation $\prod g_j^{n_j}=1$.
\end{lemma}
\begin{proof}
We give the character construction to make the dual and its fibres
explicit. Suppose a character $\chi$ is already defined on a
subgroup $B$, and choose $g\notin B$. Let $r$ be the smallest
positive integer such that $g^r\in B$. Every element of
$\langle B,g\rangle$ has a unique presentation $bg^j$ with
$b\in B$ and $0\leq j<r$: equality of two such presentations
would place $g^{j-j'}$ in $B$. For each of the $r$ distinct complex
roots $z$ of $z^r=\chi(g^r)$, define
$\widetilde\chi(bg^j)=\chi(b)z^j$.
The relation at $j=r$ proves multiplicativity, and restriction and
the value at $g$ recover $\chi,z$ uniquely. Starting at the trivial
subgroup and adjoining elements until reaching $G$ constructs
exactly $\prod r=|G|$ characters, because the subgroup order also
multiplies by $r$ at each step. To separate $u\ne1$ of order $d$,
first define the character on $\langle u\rangle$ by
$u\mapsto e^{2\pi i/d}$, then extend by this construction.
Multiplication by a character nontrivial at $u$ permutes
$\widehat G$ and multiplies the sum in (T9) by a number other
than one; hence that sum is zero. At $u=1$ every summand is one.
Finally a homomorphism is determined by its values on the retained
generators. The relation equations in (T10) are exactly what makes
the inverse formula independent of presentation; multiplication
then proves it is a character. This proves every image, inverse,
and singleton-fibre assertion.
\end{proof}

On $\widehat G^2$ the source matrix is the exact endomorphism
\[
 J_G(\chi,\psi)=(\chi^3\psi,\chi\psi^5).             \tag{T11}
\]
Under the coordinate embedding
$E_2(\chi,\psi)=((E\chi)_1,(E\psi)_1,\ldots,
(E\chi)_s,(E\psi)_s)$ it is the restriction of the $s$-block
torus map $\Phi^{\oplus s}$. Thus $E_2J_G=\Phi^{\oplus s}E_2$
coordinate by coordinate, with an explicitly proved inverse for
$E_2$ on its image. This is the precise morphism between the
continuous cover and the finite character group.

\begin{theorem}[Exact arithmetic aliasing and the original packet fibres]
\label{thm:ns-es-packet}
For a function $f$ on $\widehat G^2$, define
\[
 \widehat f(u,v)=\frac1{|G|^2}
   \sum_{\alpha,\beta\in\widehat G}
       f(\alpha,\beta)\overline{\alpha(u)\beta(v)}.
\]
Then its mean after the finite restriction (T11) is exactly
\[
 \frac1{|G|^2}\sum_{\chi,\psi}f(J_G(\chi,\psi))
   =\sum_{\substack{u\in G\\u^{14}=1}}
           \widehat f(u,u^{-3}).                     \tag{T12}
\]
Each one-coordinate marginal of $J_G$ is uniform, although their
joint distribution need not be uniform. The complete kernel is
$\{(\chi,\chi^{-3}):\chi^{14}=1\}$; a target $(\alpha,\beta)$
is in the image precisely when $\alpha^5\beta^{-1}$ is a
fourteenth power in $\widehat G$. Its complete fibre consists of
$(\chi,\alpha\chi^{-3})$ for all fourteenth roots
$\chi^{14}=\alpha^5\beta^{-1}$.

Retain an original shell $p=12h+1$ prime, $h\geq1$,
$3h+1\leq a\leq9h$, and $R=4a-p$. Put $G=U(R)$ and retain
the entire prime exponent box
\[
 {\cal B}_a=\prod_{\ell\mid a}\{0,\ldots,2v_\ell(a)\},\quad
 U_a(e)=\prod_{\ell\mid a}\ell^{e_\ell},\quad
 c_a(g)=\#\{e\in{\cal B}_a:[U_a(e)]_R=g\}.
\]
Here $R\geq3$ is odd, $a<p$, and $\gcd(a,R)=1$, since any
common prime divisor divides $p=4a-R$ and cannot divide $a<p$.
In particular every listed prime factor is a unit. Define the
actual packet function at target $g\in G$ by
\[
 F_{a,g}(\chi)=\chi(g^{-1})
       \prod_{\ell\mid a}\sum_{e=0}^{2v_\ell(a)}\chi([\ell]_R)^e.
                                                               \tag{T13}
\]
The exact one-step count is
\[
 c_a(g)=\frac1{|G|}\sum_\chi F_{a,g}(\chi)
       =\frac1{|G|^2}\sum_{\chi,\psi}F_{a,g}(\chi^3\psi).
                                                               \tag{T14}
\]
Every original exponent is retained in (T13), and the counted
fibre maps bijectively to the original divisors of $a^2$ congruent
to $g$, with inverse $u\mapsto(v_\ell(u))_{\ell\mid a}$.

For two such finite packets on this same group, with coefficients
$c_a,c_b$ and targets $g,h$, where all factors of $b$ are also
units modulo this same $R$, the exact joint mean is
\[
 \frac1{|G|^2}\sum_{\chi,\psi}
    F_{a,g}(\chi^3\psi)F_{b,h}(\chi\psi^5)
   =\sum_{\substack{u\in G\\u^{14}=1}}
          c_a(gu)c_b(hu^{-3}).                       \tag{T15}
\]
All summands are nonnegative integers. Their complete pair fibre
consists of the original exponent pairs $(e,f)$ with
$[U_a(e)]_R=gu$, $[U_b(f)]_R=hu^{-3}$, indexed by $u^{14}=1$.
The index is uniquely recovered as $u=g^{-1}[U_a(e)]_R$.

More generally write the first row of the original integer power
$J^m$ as $(A_m,B_m)$ and set $d_m=\gcd(A_m,B_m)$, with
$J^0$ having row $(1,0)$ and $d_0=1$. Then
\[
 \frac1{|G|^2}\sum_{\chi,\psi}
    F_{a,g}(\chi^{A_m}\psi^{B_m})
       =\sum_{u^{d_m}=1}c_a(gu).                     \tag{T16}
\]
The powers and their integer row entries are retained, so the
additional torsion terms are visible at every level.
\end{theorem}
\begin{proof}
Formula (T9) makes the functions $(\alpha,\beta)\mapsto
\alpha(u)\beta(v)$ an orthonormal family of $|G|^2$ elements
in the $|G|^2$-dimensional space of all functions on
$\widehat G^2$. They are therefore a basis and have the displayed
Fourier coefficients. The pullback of this basis element is
\[
 (\chi^3\psi)(u)(\chi\psi^5)(v)
       =\chi(u^3v)\psi(uv^5).
\]
By (T9) its mean is one exactly when $u^3v=uv^5=1$.
The first equation forces $v=u^{-3}$; substituting in the second
gives $u^{-14}=1$. This proves (T12). For fixed $\chi$ the first
output $\chi^3\psi$ runs bijectively through $\widehat G$ as
$\psi$ varies, so its marginal is uniform. For fixed $\psi$ the
second output $\chi\psi^5$ likewise runs bijectively as $\chi$
varies. Solving the first target equation gives
$\psi=\alpha\chi^{-3}$; substituting in the second yields
$\beta=\alpha^5\chi^{-14}$, proving the image and full fibres.
The zero-target fibre yields the stated kernel.

Expanding the finite product (T13) gives precisely
$\sum_{e\in{\cal B}_a}\chi(g^{-1}[U_a(e)]_R)$.
Applying (T9) proves its mean is $c_a(g)$; applying the first
marginal calculation proves the second equality in (T14).
Unique prime factorization proves the divisor bijection and the
displayed inverse on every fibre. In the product of the two
expanded packets, (T9) keeps exactly the exponent pairs whose
quotients $u=g^{-1}[U_a(e)]_R$ and $v=h^{-1}[U_b(f)]_R$
satisfy $u^3v=uv^5=1$. The preceding calculation gives precisely
the indexed fibres in (T15). No summand or exponent allocation is
discarded. Finally the mean of
$\chi(u)^{A_m}\psi(u)^{B_m}$ is one precisely when
$u^{A_m}=u^{B_m}=1$. These equations are equivalent to
$u^{d_m}=1$: Bezout's integer relation yields one implication,
and divisibility of both row entries by $d_m$ yields the other.
Expansion of the single packet now proves (T16).
\end{proof}

\begin{corollary}[Complete-preimage coefficient extraction and an original-shell discrepancy]
\label{cor:ns-es-cover-repair}
Let $D=E_2(\widehat G^2)\subset\mathbb T^{2s}$ and let
$H_m=\{z\in\mathbb T^{2s}:(J^m)^{\oplus s}z\in D\}$ be
the full inverse image under the block map $(J^m)^{\oplus s}$. For any function $f$ on
$\widehat G^2$, the exact identity is
\[
 \frac1{|H_m|}\sum_{z\in H_m}
   f\!\left(E_2^{-1}((J^m)^{\oplus s}z)\right)
   =\frac1{|G|^2}\sum_{\chi,\psi}f(\chi,\psi),
 \qquad |H_m|=14^{ms}|G|^2.                           \tag{T17}
\]
All lifts are given by the product of (T8), and $E_2^{-1}$ is
the coordinate formula (T10); thus every map in (T17) is typed
on its full domain. For $f(\chi,\psi)=F_{a,g}(\chi)$,
(T17) equals $c_a(g)$; for
$f(\chi,\psi)=F_{a,g}(\chi)F_{b,h}(\psi)$ it equals
$c_a(g)c_b(h)$, with no extra terms.

In the original shell $p=13$, $a=4$, $R=3$, the full box is
$e_2=0,1,2,3,4$. The residues of its five divisors
$1,2,4,8,16$ are $1,2,1,2,1$. Hence $c_a(1)=3$ and
$c_a(2)=2$, and both original targets $-4^{-1}$ and $-a$ are
$2$ in $U(3)$. Formula (T15) with both packets and both targets
equal to these values is $2\cdot2+3\cdot3=13$, whereas
the independent target count is $2\cdot2=4$.
Moreover
$J^2=\left(\begin{smallmatrix}10&8\\8&26\end{smallmatrix}\right)$,
so (T16) at $m=2$ is $c_a(2)+c_a(1)=5$, whereas the
original coefficient remains $2$. Formula (T17) returns $2$
and $4$, respectively, at every cover level.
\end{corollary}
\begin{proof}
Apply the product-cover assertion of Proposition~\ref{prop:ns-torus-powers}
to the subgroup $D$. Lemma~\ref{lem:ns-finite-dual} proves
$|D|=|G|^2$ and the precise inverse used in (T17). Every target
therefore occurs $14^{ms}$ times, which proves the identity by
fibre counting. The one-packet conclusion is (T14); the two-packet
conclusion follows by separating the independent finite sums.
The displayed residues directly give the two coefficients in
the example. Both elements of $U(3)$ satisfy $u^{14}=1$;
in (T15) their two contributions are respectively $4$ and $9$.
The first row of $J^2$ has gcd $2$, and both group elements
satisfy $u^2=1$, proving its stated mean. All original five
exponents and both nonzero coefficient fibres have been listed.
\end{proof}

The bridge therefore supplies an exact finite sampling rule for the
original shell packet and an explicit obstruction to transferring
joint averages or repeated covers on an unchanged finite torsion set.
It proves neither global shell occupancy nor a PDE assertion.
